\addcontentsline{toc}{subsubsection}{Solution 50. The Golden Staircase}

\begin{dlsummary}[title={\bfseries Solution 50. The Golden Staircase},breakable]

\begin{proof}
The unique solution is
\[
\boxed{
f(n)=\left\lfloor\varphi n\right\rfloor,
\qquad
\varphi=\frac{1+\sqrt5}{2}.
}
\]

\paragraph{1. Prove that a limiting slope exists.}

Repeated application of (1) gives
\[
m f(n)\le f(mn)\le m f(n)+m-1.
\]
Interchanging \(m,n\) and combining the bounds yields
\[
\frac{f(n)}n<\frac{f(m)}m+\frac1m,
\qquad
\frac{f(m)}m<\frac{f(n)}n+\frac1n.
\]
Thus the sequence \(f(n)/n\) is Cauchy. Write
\[
\alpha=\lim_{n\to\infty}\frac{f(n)}n.
\]
Also, superadditivity gives \(f(n)\ge n f(1)\ge n\), so \(\alpha\ge1\).

For fixed \(n\), divide
\[
k f(n)\le f(kn)\le k f(n)+k-1
\]
by \(k\), and let \(k\to\infty\). We obtain
\[
\boxed{\alpha n-1\le f(n)\le\alpha n.} \tag{3}
\]

\paragraph{2. Use composition to determine the slope.}

Divide (2) by \(n\):
\[
\left|
\frac{f(f(n))}{n}-\frac{f(n)}n-1
\right|\le\frac1n.
\]
Since \(f(n)\ge n\), its values tend to infinity, and hence
\[
\frac{f(f(n))}{n}
=
\frac{f(f(n))}{f(n)}\frac{f(n)}n
\longrightarrow\alpha^2.
\]
Therefore
\[
\alpha^2=\alpha+1.
\]
As \(\alpha\ge1\), necessarily \(\alpha=\varphi\).

\paragraph{3. Recover the exact integer values.}

By (3),
\[
\varphi n-1\le f(n)\le\varphi n.
\]
Because \(\varphi n\) is irrational, this interval contains exactly one integer:
\[
f(n)=\lfloor\varphi n\rfloor.
\]

\paragraph{4. Verify the solution.}

The elementary floor inequality
\[
\lfloor a\rfloor+\lfloor b\rfloor
\le\lfloor a+b\rfloor
\le\lfloor a\rfloor+\lfloor b\rfloor+1
\]
proves (1).

For (2), write
\[
k=\lfloor\varphi n\rfloor,
\qquad
t=\varphi n-k\in(0,1).
\]
Using \(\varphi^2=\varphi+1\),
\[
\varphi k
=n+k-(\varphi-1)t.
\]
Since \(0<(\varphi-1)t<1\),
\[
f(f(n))=\lfloor\varphi k\rfloor=n+k-1=n+f(n)-1.
\]
Thus (2) holds, with absolute difference exactly \(1\).
\end{proof}
\end{dlsummary}

\begin{takeaways}
\begin{itemize}
\item \textbf{Compare a common input:} Two estimates for \(f(mn)\) control the ratios \(f(n)/n\).
\item \textbf{Bounded errors can reveal exact slopes:} The additive error disappears after normalisation.
\item \textbf{Composition determines an algebraic number:} The limiting slope must satisfy \(\alpha^2=\alpha+1\).
\item \textbf{Integrality finishes the classification:} A narrow interval identifies each function value exactly.
\end{itemize}
\end{takeaways}
